\newpage
\section{Data and Data Audit}\label{ch:data}
\subsection{The NIFTY Benchmark}
The experiments use the NIFTY dataset \cite{saqur2024} (Hugging Face identifier \texttt{raeidsaqur/NIFTY}). The market target is the SPY exchange-traded fund. The dataset contains 2{,}111 rows with the official chronological split of Table~\ref{tab:splits}. Each row provides a text prompt containing a short history of daily market statistics (open, high, low, close, adjusted close, volume, percentage change and several technical indicators), a block of news headlines, a categorical label (Fall, Neutral or Rise) and a numeric \texttt{pct\_change}. The label rule is: Fall if the change is below $-0.5\%$, Rise if it is above $+0.5\%$, and Neutral otherwise.

\begin{table}[H]
\centering
\begin{tabular}{lrrrr}
\toprule
Split & Rows (raw) & Rows (used) & First date & Last date\\
\midrule
Train & 1{,}477 & 1{,}462 & 2010-01-06$^{*}$ & 2017-06-27\\
Validation & 317 & 317 & 2017-06-28 & 2019-02-12\\
Test & 317 & 317 & 2019-02-13 & 2020-09-21\\
\midrule
Total & 2{,}111 & 2{,}096 & & \\
\bottomrule
\end{tabular}
\caption[Official NIFTY splits and rows used]{Official NIFTY splits. Fifteen early training rows are dropped because the 20-day regime features need 20 days of history (Section~\ref{sec:windows}). $^{*}$First date of the raw rows; the dropped rows span 2010-01-06 to 2010-02-02, so the first row used is later.}\label{tab:splits}
\end{table}
\figpage{fig_dataset_timeline.png}{\textwidth}{Official chronological split and the stress-period case study window.}

\subsection{The Label Is the Next Trading Day's Return}
A first check established what the label actually measures, because the timing convention determines what may legitimately be used as an input. For each row, the \texttt{pct\_change} target was compared with several candidate returns computed from the market data embedded in the contexts. The result on 2{,}110 rows with a defined next day was unambiguous:
\begin{center}
\begin{tabular}{lc}
\toprule
Candidate series & Share within $10^{-3}$ of the target\\
\midrule
Next-day close-to-close return & \textbf{1.000}\\
Same-day close-to-close return & 0.066\\
Previous-day close-to-close return & 0.081\\
Same-day open-to-close return & 0.087\\
Context \texttt{pct\_change} of the same date & 0.064\\
\bottomrule
\end{tabular}
\end{center}
\begin{defbox}
\textbf{Timing convention.} The label of a row dated $t$ is the close-to-close return from day $t$ to the next trading day $t{+}1$. A price window may therefore end at day $t$ (the most recent information available before the predicted move). This is called $\mathrm{LAG}=0$ in the code. The benchmark's own prompt uses a window ending at $t{-}1$; this alternative ($\mathrm{LAG}=1$) is implemented as a robustness option but was not run for this report.
\end{defbox}
Ten of the 2{,}111 rows lie exactly on the $\pm0.5\%$ boundary and disagree with the label rule, most likely through rounding of \texttt{pct\_change} to four decimals (the unrounded values were not inspected; for example 2012-06-14 and 2017-11-14); they were kept. The 2{,}110 valid comparisons exclude one row, the final test date 2020-09-21, whose next trading day lies beyond the recovered series.

\subsection{Recovering a Daily Price Series}\label{sec:windows}
The dataset is not a daily series: target dates are separated by one to four business days (1{,}544 consecutive pairs, 469 gaps of two, 77 of three and 20 of four), and there are 2{,}111 rows against about 2{,}700 NYSE trading days in the same range. However, every row carries several trading days of market context (context lengths 1--8 days, 5--8 for almost all rows), and consecutive contexts overlap. Merging all contexts yields \textbf{2{,}695 distinct trading days} from 2010-01-05 to 2020-09-18. Where the same day appears in several contexts the values agree exactly (maximum disagreement in close: 0.0). Checked against the NYSE calendar, the recovered series contains every session in this range except one (2018-12-31), which is not interpolated. The fourteen other dates flagged by the first audit (Good Fridays, 2012-10-29/30 and 2018-12-05) are market closures, not missing data.

Price windows are then built from this daily series, so a ``10-day'' window means ten consecutive \emph{trading} days (the maximum calendar span of a window is 18 days), not ten rows of the sparse dataset. A window for a row dated $t$ ends at the last recovered trading day not after $t$. Because 20 days of history are needed for the regime features, the 15 earliest rows are dropped, leaving 1{,}462, 317 and 317 rows (Table~\ref{tab:splits}).

\begin{warnbox}
\textbf{Leakage guards.} The code asserts for every window and for every regime feature that its last day is not after the row's date, and that all P2/P3 splits are chronological. A further check confirms that the first column of each window equals the raw close-to-close return stored by the first processing step.
\end{warnbox}

\subsection{Class Balance Shifts Across Splits}
Table~\ref{tab:classes} shows that the label mix is not stable. The Neutral share is 53.5\% in training but 45.1\% in the test period, which contains the 2020 shock.
\begin{table}[H]
\centering
\begin{tabular}{lccc}
\toprule
Split & Fall & Neutral & Rise\\
\midrule
Train & 0.205 & 0.535 & 0.260\\
Validation & 0.167 & 0.596 & 0.237\\
Test & 0.230 & 0.451 & 0.319\\
\bottomrule
\end{tabular}
\caption{Label shares by official split.}\label{tab:classes}
\end{table}
Always predicting the training majority class (Neutral) therefore reaches an accuracy of 0.451 on the test set but a Macro-F1 of only $0.207$, because two of the three classes get an F1 of zero. For this reason Macro-F1 is the primary metric and a \emph{uniform random guesser} (Macro-F1 $\approx0.333$) is used as the meaningful floor.

\subsection{News Headlines}
Each row carries a block of daily headlines (mean about 207 headlines per row in the raw text over all 2{,}111 rows, with the number depending strongly on the era). Over the 2{,}096 modelling rows the raw volume falls sharply across the splits: the mean number of headlines per day is about 277 in the training period (maximum 1{,}627), 54 in validation and 26 in test. To bound compute and reduce this imbalance, at most 50 headlines per day are sampled at random (fixed seed 1234) before encoding; after this cap the mean is 49.0 (train), 46.4 (validation) and 26.2 (test) headlines per day. Across all splits 94{,}652 headlines were embedded and no day was left without a headline.

Mean FinBERT sentiment (probability of positive minus probability of negative) also moves across periods: $-0.068$ in training, $-0.095$ in validation and $-0.184$ in test, and in March 2020 alone it falls to $-0.319$ against $-0.176$ for the rest of the test set. As a diagnostic of how much information the daily embedding carries, the average cosine similarity (ratio of the mean dot product to the product of mean norms) between embeddings of consecutive dataset rows is 0.963, against 0.951 for random pairs of days, so the mean-pooled embedding varies only slightly from day to day. These shifts in news volume and tone are a confound for any news result and are returned to in Chapter~\ref{ch:discussion}.

\subsection{Exploratory Overview}
\figpage{eda_overview.png}{\textwidth}{Exploratory overview of the benchmark (price history, label mix and news volume by split).}
